% Integrated FF.44; review copy only, do not input separately.
\begin{claim}{FF.44}{Proved}\label{clm:FFmetabelianaffine}
Let \(E/B_p\) be finite. Every
pinned automorphism of
\(\mathbb A^1_{E,\acute et}\) over
\(BG_E\), with its relative
structural isomorphism retained,
has actual geometric point function
\[
                          x\longmapsto ux+v
                          \qquad
                          (u\in E^\times,\ v\in E).
                                                        \tag{FF.44.1}
\]
The two coefficients are unique.
This is an existence theorem for
the actual point function in the
automorphism case. It does not
assert that the entire topos
automorphism equals the one induced
by (FF.44.1), and does not classify
arbitrary pinned morphisms.
\end{claim}
\begin{proof}
Write \(f:\overline{\mathbb Q}
\to\overline{\mathbb Q}\) for
the actual point function. It is
a bijection, because the inverse
topos automorphism supplies its
inverse. It is \(G_E\)-equivariant,
and for every finite \(L/E\)
it restricts to a bijection of
\(L\), by pinned relative point
evaluation \citep[Proposition~3.6]{FF_CHW2024}.
In particular \(f(0),f(1)\in E\)
are distinct. Postcompose with
the affine scheme automorphism
\[
                   T\longmapsto
                   \frac{T-f(0)}{f(1)-f(0)}.
\]
It suffices to prove that the
new point function is the identity.
Henceforth
\[
                            f(0)=0,\qquad f(1)=1.
                                                        \tag{FF.44.2}
\]

\emph{The cohomology of each puncture.}
Let \(a\in\overline{\mathbb Q}\),
and choose a finite extension
\(L/E\) containing \(a\). Then
\(f(a)\in L\). The scalar-extended
topos automorphism restricts to an
equivalence
\[
 (\mathbb A^1_L\setminus\{a\})_{\acute et}
       \longrightarrow
 (\mathbb A^1_L\setminus\{f(a)\})_{\acute et}.
                                                        \tag{FF.44.3}
\]
Indeed its bijective geometric point
function takes precisely \(a\)
to \(f(a)\), so the inverse-image
open is precisely the displayed
source. Geometric continuous
\(p\)-adic cohomology in degree
one, with coefficients
\(\mathbb Z_p(1)\), is free of
rank one, with basis
\(\kappa(T-a)\). The equivalence
(FF.44.3) therefore acts on that
integral geometric line by a unit
\(\alpha_a\in\mathbb Z_p^\times\).
The unit is unchanged by further
finite scalar extension; equivalently
one may compute after the algebraic
scalar extension to
\(\overline{\mathbb Q}\), obtained
as the inverse limit of finite
etale base extensions.

The Hochschild--Serre sequence and
the global lift \(\kappa(T-a)\)
show that the pullback of
\(\kappa(T-f(a))\), with rational
coefficients, is
\[
                         \alpha_a\kappa(T-a)+\eta_{a,L},
           \qquad \eta_{a,L}\in H^1(G_L,\mathbb Q_p(1)).
                                                        \tag{FF.44.4}
\]
The relative structural isomorphism
identifies the coefficient system
and its base pullback here. Evaluating
at any \(x\in L\setminus\{a\}\)
gives
\[
       \kappa_L(f(x)-f(a))
          =\alpha_a\kappa_L(x-a)+\eta_{a,L}.
                                                        \tag{FF.44.5}
\]
These equations always compare points
and classes over one common finite
base \(L\).

\emph{Recovering the affine lines.}
Put \(F=K_p^{\rm ab}\). For distinct
\(a,b\) and \(x\notin\{a,b\}\)
in \(\overline{\mathbb Q}\), choose
one finite \(L/E\) containing all
three. Subtracting the evaluations
at \(x\) and \(b\) in (FF.44.5)
cancels the arithmetic term and gives
\[
 \kappa_L\!\left(\frac{f(x)-f(a)}{f(b)-f(a)}\right)
     =\alpha_a
          \kappa_L\!\left(\frac{x-a}{b-a}\right).
                                                        \tag{FF.44.6}
\]
The corresponding comparison using
the puncture \(b\) gives
\[
 \kappa_L\!\left(\frac{f(b)-f(x)}{f(b)-f(a)}\right)
     =\alpha_b
          \kappa_L\!\left(\frac{b-x}{b-a}\right).
                                                        \tag{FF.44.7}
\]
The coefficients are nonzero. By
Claim~\ref{clm:FFmetabeliankernel},
both classes on the right vanish
exactly when \((x-a)/(b-a)\in F\).
Both on the left vanish exactly
when \((f(x)-f(a))/(f(b)-f(a))\in F\).
Including the endpoints, and using
the inverse automorphism for the
reverse inclusion, we obtain
\[
                 f\bigl(a+(b-a)F\bigr)
                     =f(a)+(f(b)-f(a))F.
                                                        \tag{FF.44.8}
\]
Thus \(f\) preserves all affine
\(F\)-lines, in both directions.

\emph{An elementary affine-geometry argument.}
We give the argument also in this
infinite-dimensional setting. Regard
\(\overline{\mathbb Q}\) as a vector
space over \(F\). Its dimension
is at least two. A bijection fixing
zero and preserving lines maps any
plane spanned by independent vectors
\(x,y\) onto the plane spanned by
\(f(x),f(y)\). To see containment,
each point off the two axes is on
a line joining a point of \(Fx\)
to a point of \(Fy\); for example
\(sx+ty\) is on the line joining
\(2sx\) and \(2ty\). Use the
inverse bijection for equality.
Noncollinearity is preserved as well.
Within such a plane, parallel lines
are characterized by disjointness,
so parallelism is preserved.
The parallelogram intersections give
\[
                      f(x+y)=f(x)+f(y)
                   \quad\text{if }x,y\text{ are independent.}
                                                        \tag{FF.44.9}
\]
For dependent vectors, introduce a
vector \(z\) outside their line.
First apply (FF.44.9) to
\(z=x+(-x+z)\) and to \(-x+z\)
to obtain \(f(-x)=-f(x)\).
If \(x+y\ne0\), compute
\(f(x+y+z)\) both as
\(f(x+y)+f(z)\) and as
\(f(x)+f(y+z)\). This gives
additivity also for dependent vectors;
the zero cases follow immediately.

For each nonzero vector \(x\),
line preservation defines a bijection
\(\sigma_x:F\to F\) by
\(f(cx)=\sigma_x(c)f(x)\).
For independent \(x,y\), compare
the two additive expansions of
\(f(c(x+y))\) to see
\(\sigma_x=\sigma_y=\sigma_{x+y}\).
For dependent nonzero vectors,
use a third independent vector.
Thus a single \(\sigma\) works
for every vector. Additivity of
\(f\) makes \(\sigma\) additive,
and comparing \(f(cd x)\) in
two ways makes it multiplicative.
It is bijective and takes one
to one, hence is an automorphism
of \(F\). In view of (FF.44.2),
\[
      f(x+y)=f(x)+f(y),\qquad
      f(cx)=\sigma(c)f(x)\ (c\in F),\qquad
      f|_F=\sigma.
                                                        \tag{FF.44.10}
\]
In particular \(f(F)=F\) and
\(f\) fixes \(\mathbb Q\).
Only two-dimensional spans were
used, so there is no finite-dimension
assumption hidden in this argument.

\emph{Recovering multiplication.}
In (FF.44.5) take \(a=0\).
Evaluating at \(x=1\) and using
(FF.44.2) gives \(\eta_{0,L}=0\).
Consequently, on every common finite
extension \(L/E\),
\[
                            \kappa_L(f(x))
                                   =\alpha_0\kappa_L(x).
                                                        \tag{FF.44.11}
\]
Let \(x,y\ne0\), and choose
one such \(L\) containing both.
It also contains all their images
under \(f\). For each rational
\(c\) with \(x+c\ne0\),
multiplicativity of Kummer classes
and (FF.44.11) yield
\[
       \kappa_L(f((x+c)y))
                     =\kappa_L(f(x+c))+\kappa_L(f(y)).
\]
Using (FF.44.10), this becomes
\[
       \kappa_L\!\left(\frac{f(xy)}{f(y)}+c\right)
                     =\kappa_L(f(x)+c).
                                                        \tag{FF.44.12}
\]
The equalities hold for all rational
translates at which both sides
are defined. The final assertion
of Claim~\ref{clm:FFmetabeliankernel}
therefore gives either
\[
                        \frac{f(xy)}{f(y)}=f(x)
\]
or both displayed values lie in
\(F\). If \(x\notin F\),
then \(f(x)\notin F\), so
only the first alternative is
possible. If \(x\in F\),
semilinearity already gives
\(f(xy)=\sigma(x)f(y)=f(x)f(y)\).
Zero factors are harmless. Hence
\(f\) is multiplicative everywhere.

The normalized point function is
therefore a field automorphism of
\(\overline{\mathbb Q}\).
Its \(G_E\)-equivariance says exactly
that it centralizes \(G_E\).
Claim~\ref{clm:FFmetabeliantest}
applies to every finite \(E/B_p\)
and makes this centralizer trivial.
Thus the normalized \(f\) is
the identity. Undoing the affine
normalization gives (FF.44.1), with
\(u=f_{
m original}(1)-f_{
m original}(0)\)
and \(v=f_{
m original}(0)\).
\end{proof}

The unresolved automorphism kernel
now has a precise form: after
composition with a unique affine
scheme automorphism, an additional
pinned automorphism, if one exists,
must fix every geometric point.
The proof above does not identify
its action on the entire etale
site, and does not provide the
missing existence theorem for
noninvertible pinned maps.
